\documentclass[10pt,a4paper]{amsart}
\usepackage[margin=19mm]{geometry}
\usepackage{amsmath,amssymb,mathtools}
\usepackage[scr=rsfs,cal=euler]{mathalfa}
\usepackage{xfrac}
\usepackage[unicode,hidelinks,bookmarks=false]{hyperref}
\newcommand{\ie}{i.e.\ }
\newcommand{\cf}{cf.\ }
\newcommand{\resp}{respectively\ }
\newcommand{\Subsection}{\subsection}
\providecommand{\nobreakdash}{\nobreak\hbox{-}}
\setlength{\emergencystretch}{2em}
\multlinegap0pt
\usepackage{bm}
\usepackage{bbm}
\usepackage{dsfont}
\usepackage{xspace}
\usepackage{cancel}
\usepackage{mathtools}
\usepackage{amsthm}
\makeatletter
\@ifundefined{thm}{\newtheorem{thm}{Thm}}{}
\@ifundefined{rem}{\newtheorem{rem}[thm]{Remark}}{}
\makeatother
\makeatletter
\newcommand{\Cay}[2]{\operatorname{Cay}\left(#1,#2\right)}
\expandafter\ifx\csname claimproof\endcsname\relax
\newenvironment{claimproof}{\renewcommand{\qedsymbol}{$\blacksquare$}\begin{proof}}{\end{proof}}
\fi
\makeatother
\title[Short hierarchically hyperbolic groups I: uncountably many c]{Short hierarchically hyperbolic groups I: uncountably many coarse median structures}
\author{Giorgio Mangioni}
\date{Source: arXiv:2410.09232v5 (CC BY 4.0)}
\begin{document}
\maketitle

\noindent\emph{Source and licence.} Short hierarchically hyperbolic groups I: uncountably many coarse median structures. Version arXiv:2410.09232v5 (CC BY 4.0), distributed under the Creative Commons Attribution 4.0 International licence, as recorded in the arXiv licence metadata for this exact version and shown on the abstract page https://arxiv.org/abs/2410.09232v5. Full rights notice: licenses/arxiv-2410.09232-CC-BY-4.0.txt.\par\smallskip

\noindent\textit{Editorial scope.} Counted unit: the Rem whose source label is \texttt{rem:mG\_trivial\_on\_adjacent\_Zw} in the article (source0.tex lines 1000--1002, source label \texttt{rem:mG\_trivial\_on\_adjacent\_Zw}), delivered together with the article's own complete original proof of it (source0.tex lines 1004--1025, one \texttt{proof} environment of 22 lines). No mathematics is written, repaired or re-derived: statement and proof are the article's own text line for line. Required context retained uncounted: none. Every definition, display and label of those lines is retained unchanged. Omitted from this package and not redistributed: 1--999, 1003--1003, 1026--2047. Nothing inside a retained excerpt is omitted or altered: every retained body is delivered exactly as the article's lines are written, so no omission receipt of any kind accompanies this package. Citations: the retained excerpts carry no citation command at all, so no bibliography entry of the article is transcribed and no reference is redistributed. Reference adaptation: none was needed and none was made; every reference inside the delivered unit resolves inside it, so no unresolved reference or citation is produced. Printed slips kept literal: no printed word, symbol, hypothesis or number of the counted statement or of its proof was altered. Layout and class: the article's own class is \texttt{amsart} with packages including babel, inputenc, svg, hyperref, tabularx, ifthen, amssymb, amsfonts, amsmath, amsthm, amscd, amsxtra, latexsym, mathabx; the wrapper is the lane's pinned \texttt{amsart} editorial wrapper at 10pt on A4, which restates the theorem-like environments the retained text needs and adds the article's own macro definitions that the retained text needs (\texttt{\textbackslash Cay}), with extra wrapper packages (bm, bbm, dsfont, xspace, cancel, mathtools). The delivered numbering is the unit's own. Assets: no retained span contains a figure environment, a graphics-inclusion command, a PDF-page inclusion, a table or a TikZ picture; no image file is copied into this package, the package declares no asset dependency and no image file is redistributed with it. Licence: version arXiv:2410.09232v5 of this article is distributed under the Creative Commons Attribution 4.0 International licence, as recorded in the arXiv licence metadata for this exact version and shown on the abstract page https://arxiv.org/abs/2410.09232v5. A full rights notice is delivered at licenses/arxiv-2410.09232-CC-BY-4.0.txt. Characters: every retained excerpt is pure ASCII, so the delivered engine, XeLaTeX with its default Latin Modern text font, reports no missing character.
\begin{rem}\label{rem:mG_trivial_on_adjacent_Zw}
    Later we shall also need that, if $e\in E$ is such that $m(geg^{-1})=0$ for every $g\in G$, then clearly $m^G(e)=0$. 
\end{rem}

\begin{proof}
Firstly, each $m(g_i \cdot g_i^{-1})$ is a homogeneous quasimorphism, being the composition of a group automorphism and $m$. Then $m^G$ is a homogeneous quasimorphism, since it is defined as a linear combination of homogeneous quasimorphisms. Furthermore, by construction $m^G(z)=m(z)$, so $m^G$ is itself unbounded on $\langle z\rangle$.


Now we claim that the absolute value of $m^G$ is $G$-invariant. Indeed, for every $g\in G$ and every $h\in E$, we have that
$$m^G(g h g^{-1})=\sum_{i=1}^k\varepsilon(g_i) m(g_ighg^{-1}g_i^{-1}).$$
Let $g_ig=g_{j(i)} h_{j(i)}$, for some $g_{j(i)}$ in the set of representatives and some $h_{j(i)}\in E$. Notice that $i\to j(i)$ is a permutation of the set $\{1, \ldots, k\}$, as $g_ig$ and $g_jg$ cannot lie in the same coset of $E$. Furthermore $\varepsilon(g_i)=\varepsilon(g)\varepsilon(g_{j(i)})$, as $h_{j(i)}$ commutes with $z$. Hence
$$m^G(g h g^{-1})=\varepsilon(g)\sum_{i=1}^k\varepsilon(g_{j(i)}) m(g_{j(i)}h_{j(i)}hh_{j(i)}^{-1}g^{-1}g_{j(i)})=$$
$$=\varepsilon(g)\sum_{i=1}^k\varepsilon(g_{j(i)}) m\left((g_{j(i)}h_{j(i)}g_{j(i)}^{-1})(g_{j(i)}hg_{j(i)}^{-1})(g_{j(i)}h_{j(i)}^{-1}g_{j(i)}^{-1})\right).$$
Now $g_{j(i)}h_{j(i)}^{-1}g_{j(i)}^{-1}\in E$, as $E$ is normal in $G$, and using that $m$ is invariant under conjugation by elements of $E$ we get that
$$m^G(g h g^{-1})=\varepsilon(g)\sum_{i=1}^k\varepsilon(g_{j(i)}) m\left(g_{j(i)}hg_{j(i)}^{-1}\right)=\varepsilon(g)m^G(h).$$



Now choose $C$ as above. Clearly $i_m$ is $1$-Lipschitz and coarsely surjective, so in order to show that it is a quasi-isometry we shall prove that there exists a coarsely Lipschitz retraction 
$$r_m\colon \Cay{G}{\tau_{\{m^G,C\}}\cup\{g_1,\ldots, g_k\}}\to \Cay{E}{\tau_{\{m^G,C\}}}.$$
For every $g\in G$, we can write it uniquely as $g=hg_i$ for some $h\in E$ and $g_i$ in the chosen collection of coset representatives. Then we set $r_m(g)=h$, which is the identity on $E$. 

Now suppose $g, g'\in G$ are adjacent in $\Cay{G}{\tau_{\{m^G,C\}}\cup\{g_1,\ldots, g_k\}}$. We can write $g=hg_i$ as above. If $g'=gg_j$ for some coset representative $g_j$ then, setting $g_ig_j=h_{ij}g_{l}$ for some $h_{ij}\in E$ and $g_j$ in the set of representatives, we have that $r_m(g)=h$ and $r_m(g')=hh_{ij}$ differ by $h_{ij}$, chosen in a finite subset of $E$, and in particular their distance is uniformly bounded.

If instead $g'=gh'$ for some $h'\in \tau_{\{m^G,C\}}$, then $g'=h(g_i h' g_i^{-1}) g_i$, so $r_m(g)$ and $r_m(g')$ differ by $g_i h' g_i^{-1}$, which still belongs to $\tau_{\{m^G,C\}}$ as the absolute value of $m^G$ is $G$-invariant.
\end{proof}

\end{document}
